\documentclass[11pt]{article}
\usepackage[margin=1in]{geometry}
\usepackage{amsmath,amssymb,amsthm,booktabs,hyperref}
\newtheorem{theorem}{Theorem}
\newtheorem{lemma}{Lemma}
\newcommand{\ip}[2]{\langle #1,#2\rangle}
\title{A 197580-point kissing configuration in dimension 25}
\author{}
\date{26 September 2026}
\begin{document}
\maketitle

\noindent
This note describes a finite modification of the 197579-point configuration
supplied in the public
\href{https://github.com/alexlegeartis/KissingNumbers/tree/cf14c5ef4db6059bbf2e570e3f0ce24edfda243e/verifications/improved/dim25-lens-heads}
{\texttt{KissingNumbers/dim25-lens-heads}} package.
The original configuration, its 1016 heads, and the Leech-shell construction
are inputs credited to that package. The modification below moves 552
retained equatorial points, adjoins one point, and replaces two individual
cap points.

\paragraph{Normalization and input.}
All points have squared norm 4; the kissing condition is an inner product
at most 2 between distinct points. Write the input configuration as
\[
 C_0=\{(z,0):z\in L\setminus D\}
 \cup\{(x_i,\pm1):0\le i<1016\}
 \cup\{P_0,N,S\},
\]
where $L$ is the 196560-vector minimal shell of the Leech lattice,
$|D|=1016$, $P_0=(P,0)$, and $N=(0,2)$, $S=(0,-2)$.
The head indices refer to rows of the source \texttt{heads\_X.npy}; the
mathematical values of its 44 rational heads are preserved as exact
rational numbers in \texttt{baseline-heads.json}.

Let $v$ be the first complete block's norm-6 lean and put $n=v/\sqrt6=-P/2$.
An exact integer description of this direction is
\[
v=\frac{1}{\sqrt8}(-5,-1,-1,-1,-1,1,-1,1,-1,-1,-1,1,
1,-1,-1,-1,1,1,-1,1,-1,-1,1,-1).
\]
The complete block is $U=\{u\in L:\ip uv=-3\}$, of size 552.
Every $u\in U$ belongs to $D$. For each such $u$, the point $w=u+v$ is
minimal, lies outside $D$, and satisfies $\ip uv=-3$ and $\ip uw=1$.
The corresponding old cap points have horizontal coordinate
$u+t v$, with $t=(3-\sqrt3)/6$.
All these input identities are checked by the shell and owner data.

\paragraph{The modification.}
For $u\in U$ define
\[
r=u+\sqrt{3/2}\,n=w-\sqrt{3/2}\,n,\qquad
a=\frac{\sqrt3}{2}+\frac{\sqrt2}{4},\qquad
b=\frac{\sqrt6-2}{4}.
\]
Then $r\perp n$, $\|r\|^2=5/2$, and the old cap points over $u$ are
$(r-n/\sqrt2,\pm1)$. Replace every $(w,0)$ by
\[
 W(r)=(r+a n,b),
\]
and adjoin
\[
 Q=(\sqrt3 n,-1)=(V/4,-1),
\]
where $V=\sqrt8\,v$ is the integer vector printed above. In particular,
the new point $Q$ is rational and has squared norm $48/16+1=4$.
Finally replace the upper point $(x_{984},1)$ and the lower point
$(x_{1008},-1)$ by $R_+$ and $R_-$, respectively. The certificate specifies
two integer vectors and defines
\[
 R_+=\frac{2k}{\sqrt{1007176}},\qquad
 R_-=\frac{2\ell}{\sqrt{10209}},
\]
where the following integer vectors have squared norms 1007176 and 10209:
\begin{align*}
k={}&(-336,-9,-298,286,-17,15,-39,22,277,19,364,304,9,\!\\
 &\phantom{(}-2,-24,-39,4,304,-9,-4,271,-29,9,-2,500),\\
\ell={}&(-25,-31,-33,4,3,2,36,33,3,-1,-2,-3,-4,\!\\
 &\phantom{(}-1,4,-2,1,2,4,-31,-1,-28,-3,-28,-50).
\end{align*}
The opposite old cap points $(x_{984},-1)$ and $(x_{1008},1)$ are retained.
The certificate, rather than an optimization routine or floating-point
array, defines the selected integer vectors.

\begin{lemma}[Norms and pairs not involving other cap points]
All the $W(r)$ and $Q$ have squared norm 4. They are mutually compatible,
and are compatible with all unchanged equatorial points, both poles, and
$P_0$. Each $W(r)$ is also compatible with the old cap points belonging
to the complete first block.
\end{lemma}

\begin{proof}
The identities
\[
a^2+b^2=\frac32,\qquad \sqrt3a-b=2,\qquad
-\frac a{\sqrt2}+b=-\frac34
\]
give the norm assertions and $\ip{W(r)}Q=2$.
For two different block points,
\[
\ip{W(r_i)}{W(r_j)}=\ip{r_i}{r_j}+\frac32=\ip{w_i}{w_j}\le2.
\]
Since $0<a<\sqrt{3/2}$, the horizontal component of $W(r)$ is
\[
r+a n=(1-\lambda)w+\lambda u,\qquad
\lambda=\frac{\sqrt{3/2}-a}{2\sqrt{3/2}}\in(0,1).
\]
An unchanged equatorial point $z$ is different from both $u$ and $w$.
Both $\ip uz$ and $\ip wz$ are at most 2, so their convex combination is
at most 2.

The shell has integral $\ip vz\in\{-3,-2,-1,0,1,2,3\}$.
The negative extreme layer is deleted and the positive extreme layer
is exactly the moved set. Thus every unchanged equatorial point obeys
\[
 \ip Q{(z,0)}=\frac{\ip vz}{\sqrt2}\le\sqrt2<2.
\]
These layer statements can also be checked directly on the finite integer
shell, without invoking a classification of holes or a numerical optimizer.

The remaining pure-point products are
\[
\ip{W(r)}{P_0}=-2a,\quad
\ip{W(r)}N=2b,\quad \ip{W(r)}S=-2b,
\]
\[
\ip Q{P_0}=-2\sqrt3,\quad \ip QN=-2,\quad \ip QS=2.
\]
All satisfy the required inequality. For the same block index, the
products with the old upper and lower cap points are, respectively,
\[
\frac52-\frac a{\sqrt2}+b=\frac74,
\qquad
\frac52-\frac a{\sqrt2}-b=\frac{11}{4}-\frac{\sqrt6}{2}<2.
\]
For different block indices, $\ip{r_i}{r_j}\le1/2$, so both products
are smaller still.
\end{proof}

\paragraph{The finite certificate.}
The previous lemma handles every equality created by the modification.
The remaining products to certify are:
\begin{enumerate}
\item the 552 moved points against the 2030 unchanged cap points;
\item $Q$ against those unchanged cap points;
\item $R_+$ and $R_-$ against every point retained in the final configuration,
      including one another, with their own diagonal entries omitted.
\end{enumerate}
Their observed maxima are approximately
\[
1.99937195451151,\quad 1.96592582628907,\quad
1.99472127475566,\quad 1.98752106512801.
\]
These decimals are diagnostic values, not rounded-up proof bounds.
The accompanying integer/interval certificate has checked these groups with
strict upper bounds below 2 at 192-bit Arb precision and checked both
normalizations exactly. In particular, the first two upper bounds are
strictly below 1.999372 and 1.965926.
No conclusion depends on an optimizer success flag. All pairs involving
only unchanged points are inherited from the verified source configuration.

\begin{theorem}
The specified configuration consists of 197580 distinct points on
the radius-2 sphere in $\mathbb R^{25}$ with pairwise inner products at
most 2. Consequently the kissing number in dimension 25 is at least 197580.
\end{theorem}

\begin{proof}
The lemma, the verified finite checks, and the independently checked source
certificate exhaust all pairs. The baseline checks include every head--shell
and same-cap head--head inequality, with integer or rational acceptance
decisions. Compatibility within the Leech shell follows from its lattice
minimum.
All squared norms are exactly 4, including those of $R_+$ and $R_-$ by
definition. The 552 moved points replace 552 existing points and the two
repairs replace two existing cap points; only $Q$ increases the inventory.
Thus
\[
 |C|=197579-552-2+552+2+1=197580.
\]
No duplicate can occur: two equal norm-4 vectors would have inner product
4, contradicting the verified bound 2 for their distinct indices.
\end{proof}

\end{document}
